% Clase 7 — la relacion de recurrencia salta de 2 en 2: pares e impares forman
% cadenas independientes, y cuando K = n(n+1) la cadena se corta y queda un polinomio.
\begin{tikzpicture}[scale=1,
  coef/.style={draw=figblue, fill=figblue!8, rounded corners=2pt,
               minimum width=0.86cm, minimum height=0.56cm, font=\small},
  coefcero/.style={draw=figgray!60, fill=figgray!12, rounded corners=2pt,
               minimum width=0.86cm, minimum height=0.56cm, font=\small, text=figgray}]

% ---- cadena par ----
\node[coef] (a0) at (0,0.9)   {$a_0$};
\node[coef] (a2) at (2.0,0.9) {$a_2$};
\node[coef] (a4) at (4.0,0.9) {$a_4$};
\node[coefcero] (a6) at (6.0,0.9) {$a_6$};
\node[coefcero] (a8) at (8.0,0.9) {$a_8$};
\draw[figvec] (a0) -- (a2);
\draw[figvec] (a2) -- (a4);
\draw[figvec] (a4) -- (a6);
\draw[figvec] (a6) -- (a8);
\node[figlbl] at (-1.15,0.9) {pares};

% ---- cadena impar ----
\node[coefcero] (a1) at (0,-0.9)   {$a_1$};
\node[coefcero] (a3) at (2.0,-0.9) {$a_3$};
\node[coefcero] (a5) at (4.0,-0.9) {$a_5$};
\node[coefcero] (a7) at (6.0,-0.9) {$a_7$};
\node[coefcero] (a9) at (8.0,-0.9) {$a_9$};
\draw[figvecaux] (a1) -- (a3);
\draw[figvecaux] (a3) -- (a5);
\draw[figvecaux] (a5) -- (a7);
\draw[figvecaux] (a7) -- (a9);
\node[figlbl, text=figgray] at (-1.15,-0.9) {impares};

% el corte
\draw[figred, line width=1.1pt] (5.0,1.72) -- (5.0,0.08);
\node[figlbl, figred, align=center] at (5.0,2.25) {$K=n(n+1)$\\ corta ac\'a};

\node[figlblaux, align=center] at (4.0,-1.75)
  {Con $a_1=0$ toda la cadena impar es nula: se elige una paridad y se trabaja en ella.};

\node[figlbl, align=center] at (4.0,-2.7)
  {$a_{m+2}=a_m\,\dfrac{m(m+1)-K}{(m+1)(m+2)}$\ : el salto de dos deja dos cadenas que\\[0.3em]
   no se comunican, y el numerador se anula justo en $m=n$ cuando $K=n(n+1)$.};

\end{tikzpicture}
